% !TEX root = ../main.tex
%
% CHL Borcherds products, paramodular Fricke covariance, and the
% mirror interpretation at levels N in {1,2,3,4,6}.
%
\providecommand{\Z}{\mathbb{Z}}
\providecommand{\Q}{\mathbb{Q}}
\providecommand{\cO}{\mathcal{O}}
\providecommand{\kBKM}{\kappa_{\mathrm{BKM}}}
\providecommand{\kch}{\kappa_{\mathrm{ch}}}
\providecommand{\kcat}{\kappa_{\mathrm{cat}}}
\providecommand{\kfib}{\kappa_{\mathrm{fiber}}}
\providecommand{\Sp}{\mathrm{Sp}}
\providecommand{\Fuk}{\mathrm{Fuk}}
\providecommand{\PhiCHL}{\Phi^{\mathrm{CHL}}}
\providecommand{\Dfive}{D_5}
\providecommand{\Htwo}{\mathfrak{H}_2}

\section*{CHL Fricke Covariance and the Borcherds Weight}
\label{wn:sec:chl-fricke-covariance}

\paragraph{The diagonal CHL threefold.}
Let \(S\) be a projective K3 surface, let \(E\) be an elliptic curve,
and let \(N\in\{1,2,3,4,6\}\).  Fix a symplectic automorphism
\(g_N\in \operatorname{Aut}(S)\) of order \(N\) and translation by an
\(N\)-torsion point \(t_N\in E[N]\).  The diagonal action
\[
  G_N=\langle g_N\times t_N\rangle
\]
on \(S\times E\) is free: for \(a\ne 0\), the translation \(t_N^a\)
has no fixed point on \(E\).  Thus
\[
  X_N=(S\times E)/G_N
\]
is a smooth projective threefold, and \(S\times E\to X_N\) is finite
etale of degree \(N\).  The holomorphic form \(\Omega_S\wedge dz_E\)
is \(G_N\)-invariant, since \(g_N\) is symplectic and translation fixes
\(dz_E\).  It descends to a nowhere-vanishing holomorphic three-form on
\(X_N\).

The same finite-etale calculation fixes the Hodge-column invariants:
\[
  H^q(X_N,\cO_{X_N})
  =H^q(S\times E,\cO_{S\times E})^{G_N}.
\]
The symplectic K3 automorphism is trivial on \(H^{0,0}(S)\) and
\(H^{0,2}(S)\), and the elliptic translation is trivial on
\(H^0(E,\cO_E)\) and \(H^1(E,\cO_E)\).  Hence
\[
  (h^{0,0},h^{0,1},h^{0,2},h^{0,3})(X_N)=(1,1,1,1),
\]
and therefore
\[
  \kcat(X_N)=\chi(\cO_{X_N})=0,\qquad
  \kch^{\mathrm{comp}}(X_N)=\sum_q(-1)^q h^{0,q}(X_N)=0.
\]
At \(N=1\) this is the total-space value
\(\kcat(K3\times E)=\chi(\cO_{K3})\chi(\cO_E)=2\cdot0=0\).
The fibre value \(\chi(\cO_{K3})=2\) belongs to the K3 fibre.

\paragraph{The primitive CHL Borcherds constants.}
Let \(\PhiCHL_N\) denote the primitive CHL Borcherds denominator
attached to the \(g_N\)-twined K3 Jacobi input in the
Gritsenko--Nikulin normalization.  Let \(c_N(0)\) be the zero-component
constant term of that Borcherds input.  Borcherds' singular-theta
weight theorem, in the form used in Gritsenko--Nikulin and Gritsenko's
CHL denominator construction, gives
\[
  \kBKM(\PhiCHL_N)=\operatorname{wt}(\PhiCHL_N)=\frac{c_N(0)}{2}.
\]
The primitive CHL ladder is
\[
\begin{array}{c|ccccc}
N & 1 & 2 & 3 & 4 & 6\\ \hline
c_N(0) & 10 & 8 & 6 & 4 & 2\\
\operatorname{wt}(\PhiCHL_N)=\kBKM(\PhiCHL_N) & 5 & 4 & 3 & 2 & 1
\end{array}
\]
(Borcherds 1998, Theorem~13.3; Gritsenko--Nikulin,
\emph{Automorphic forms and Lorentzian Kac--Moody algebras II},
Theorem~2.1; Gritsenko 1999, Theorem~1.12).

At \(N=1\), the input is the Eichler--Zagier weak Jacobi form
\[
  \phi_{0,1}(\tau,z)=\zeta^{-1}+10+\zeta+O(q),
\]
so \(c_1(0)=10\) and \(\kBKM(\PhiCHL_1)=5\).  The monic
Gritsenko--Nikulin product is
\[
  \Dfive(Z)=64^{-1}\Delta_5(Z),
\]
with the same weight \(5\).  The Oberdieck--Pandharipande scalar square
uses \(\Dfive^2=4096^{-1}\Delta_5^2\); the scalar \(64\) and the square
leave the primitive Borcherds characteristic of \(\Delta_5\) unchanged.

\paragraph{The Gritsenko--Clery diagonal-divisor catalogue.}
The Gritsenko--Clery simplest-divisor catalogue is a separate
eight-row classification surface.  Its row weights are
\[
  (5,2,3,1,2,\tfrac12,\tfrac32,1),
\]
with twice-weight tuple
\[
  (10,4,6,2,4,1,3,2).
\]
The two tables overlap at the common row \(F_{1,1}=\Delta_5\).  A
componentwise identification of the CHL ladder with the diagonal-divisor
catalogue requires an explicit row map: the pair \((g,h)\), the Jacobi
input, the constant term, the paramodular group, and the multiplier
system.  After those data are supplied, the same Borcherds formula
\(\kBKM=c(0)/2\) applies to that row.

\paragraph{Fricke covariance.}
Let \(K(N)\subset \Sp_4(\Q)\) be the paramodular group and let
\[
  w_N(Z)=-(NZ)^{-1},\qquad Z\in\Htwo,
\]
be the Fricke involution.  Suppose the primitive CHL denominator
\(\PhiCHL_N\) extends to the Fricke normalizer with multiplier
\(\nu_N\).  Then its Fricke transformation law is
\begin{equation}
  \PhiCHL_N(w_N Z)
  =
  \nu_N(w_N)\,\det(NZ)^{\kBKM(\PhiCHL_N)}\,\PhiCHL_N(Z).
  \label{wn:eq:chl-fricke-character}
\end{equation}
This is the definition of a paramodular form of weight
\(\kBKM(\PhiCHL_N)\) and multiplier \(\nu_N\) on the extended group,
evaluated at the Fricke element.  The exponent is the Borcherds weight
\(c_N(0)/2\).  The scalar \(\nu_N(w_N)\) is multiplier-system data.  Its
computation requires the Fricke character in the chosen automorphic
normalization; the table above fixes the weight.

For \(N=1\), \(\Delta_5\) has the Maass character
\(\nu_{\Delta_5}\) on \(\Sp_4(\Z)\), and
\[
  \Delta_5(w_1 Z)=\nu_{\Delta_5}(w_1)\det(Z)^5\Delta_5(Z).
\]
Squaring gives the trivial-multiplier Igusa form
\(\Delta_5^2=\Phi_{10}^{\mathrm{un}}\) of weight \(10\).  The reduced
Oberdieck--Pandharipande scalar transforms with inverse weight
\(-10\), because it is proportional to \(\Delta_5^{-2}\).

\paragraph{Borcherds denominator and lattice reflection.}
On the level-one lattice \(\Lambda^{2,1}_{II}\), the Weyl--Kac--
Borcherds denominator identity is
\[
  64^{-1}\Delta_5(2Z)
  =
  e^{-2\pi i(\rho,Z)}
  \prod_{\alpha\in \Delta_+}
  \bigl(1-e^{-2\pi i(\alpha,Z)}\bigr)^{\operatorname{mult}\alpha},
\]
with multiplicities read from the Fourier coefficients \(f(nm,\ell)\)
of \(\phi_{0,1}\).  The Fricke operation exchanges the two primitive
isotropic cusp directions in the orthogonal model \(O(2,3)\).  The
automorphic factor in \eqref{wn:eq:chl-fricke-character} is therefore
the line-bundle factor of this cusp exchange.

For \(N>1\), the same sentence is a typed statement only after the
twined Jacobi input, the reflection lattice, the Weyl vector, and the
multiplier have been named.  The CHL constants above determine the
weight; the Fricke eigenvalue and lattice normal form are additional
automorphic data.

\paragraph{Mirror interpretation.}
The unconditional theorem in this note is the automorphic
transformation law \eqref{wn:eq:chl-fricke-character}.  A homological
mirror statement has the following precise input.  Assume an equivariant
HMS equivalence
\[
  D^b_{G_N}\operatorname{Coh}(S\times E)
  \simeq
  D^\pi\Fuk(\check X_N)
\]
compatible with the finite etale quotient and with the lattice-polarized
period map carrying the two primitive cusp directions to the mirror
large-volume and large-complex-structure limits.  Under this hypothesis,
the Fricke involution is the automorphic avatar of the mirror exchange
of those cusp directions, and the cocycle in
\eqref{wn:eq:chl-fricke-character} is the automorphic-line factor
transported by the equivalence.

The free diagonal quotient is essential here.  The quotient map is finite
etale, so the local geometry of \(X_N\) is smooth descent from
\(S\times E\).  Fixed points of \(g_N\) on \(S\) are removed from the
quotient geometry by the free elliptic translation.  The mirror problem
is equivariant finite-etale descent plus lattice polarization.

\paragraph{Scope of the five levels.}
The present CHL statement is indexed by \(N\in\{1,2,3,4,6\}\) and uses
the primitive CHL tuple
\[
  (c_N(0))=(10,8,6,4,2),\qquad
  \bigl(\kBKM(\PhiCHL_N)\bigr)=(5,4,3,2,1).
\]
The Gritsenko--Clery diagonal-divisor rows with twice weights
\((10,4,6,2,4,1,3,2)\) form a different catalogue.  Levels outside the
primitive CHL set require their own automorphic group, multiplier, input
Jacobi form, and geometric host before a mirror-covariance statement is
available.

\paragraph{Separated invariants.}
For the \(K3\times E\) anchor the five relevant numbers are
\[
  \kcat(K3\times E)=0,\qquad
  \kch^{\mathrm{comp}}(K3\times E)=0,\qquad
  \kch^{\mathrm{Heis}}(K3\times E)=3,\qquad
  \kBKM(\Delta_5)=5,\qquad
  \kfib(K3)=24.
\]
They arise from distinct constructions: total-space structure-sheaf
Euler characteristic, compact Hodge supertrace, Heisenberg--Mukai
specialisation, Borcherds denominator weight, and the Mukai-lattice
rank of the K3 fibre.  At \(d=3\), the native CY-to-chiral output is
\(E_1\)-chiral; any \(E_2\)-braiding belongs to the derived centre,
double, or representation category after the relevant pairing and
completion data have been supplied.
